\begin{center}\begin{tikzpicture}[font=\small,>=stealth,line join=round]
\draw[gray!40,rounded corners=3pt] (0,0) rectangle (14.2,10.2);\node at (7.1,9.6) {例4：同一中心的距离差，远离中心时趋近零};
\node[blue] at (7.1,8.65) {$\displaystyle f(x)=\sqrt{(x-1)^2+9}-\sqrt{(x-1)^2+1}$};\begin{scope}[shift={(4,3.0)},x=0.5142857142857143cm,y=1.3cm]\draw[->,gray] (-6,0)--(8.5,0) node[right]{$x$};\draw[->,gray] (0,-0.6)--(0,3.4000000000000004) node[left]{$y$};\draw[gray,dashed] (-6,0)--(8,0);\node[left=3pt,fill=white,inner sep=1pt] at (0,0) {$0$};\draw[blue,very thick,domain=-6:8,samples=240] plot(\x,{sqrt((\x-1)^2+9)-sqrt((\x-1)^2+1)});\draw[gray,dotted] (1,0)--(1,2);\fill[blue] (1,2) circle(2.3pt);\node[below=4pt] at (1,0){$1$};\node[blue,above=9pt] at (1,2) {$(1,2)$};\end{scope}\node[anchor=west,align=left] at (9,4.5) {$m=n=1,\ a=3>b=1$\\[12pt]对每个实数 $x$：$f(x)>0$\\[12pt]值域 $(0,2]$\\[12pt]若两根号完全相同\\[9pt]则恒为 $0$，值域 $\{0\}$。};
\node[text width=12.8cm,align=center] at (7.1,.65) {蓝线表示题给函数；虚线 $y=0$ 是其渐近高度。相同根号的差才恒为零。};
\end{tikzpicture}\end{center}
